\documentclass[12pt]{article}
\usepackage[utf8]{inputenc}
\usepackage[T1]{fontenc}
\usepackage[margin=1in]{geometry}
\usepackage{setspace}
\usepackage{xcolor}
\usepackage[colorlinks=true, urlcolor=blue!60!black, linkcolor=blue!60!black]{hyperref}
\setlength{\parindent}{2em}
\setlength{\parskip}{1.25em}
\renewcommand{\baselinestretch}{1.0}
\begin{document}
\begin{center}
{\LARGE \textit{Employment Division v. Smith}, 494 U.S. 872 (1990)\par}
\vspace{0.8em}
{\large Case Brief --- Four-Page Edition\par}
\end{center}
\vspace{1em}

\section{Memory Jogger}

The Free Exercise Clause does not exempt individuals from complying with valid, neutral laws of general applicability, even when those laws incidentally burden central religious practices such as the sacramental use of peyote (\href{https://supreme.justia.com/cases/federal/us/494/872/}{494 U.S. 872}).

\section{Detailed Case Facts}

Alfred Smith and Galen Black were drug and alcohol rehabilitation counselors at ADAPT, a private Oregon nonprofit with a strict drug-free policy. Both were members of the Native American Church, and their sincerity was never disputed. In 1983 and 1984 each ingested a small quantity of peyote for sacramental purposes during a church ceremony, and ADAPT fired them.

Oregon denied their unemployment claims because they had been discharged for work-related ``misconduct.'' Oregon never prosecuted them, but state law made knowing possession of a Schedule I substance, including peyote, a Class B felony. The case arose during the ``War on Drugs,'' against a background of federal accommodation of Native American Church peyote use (a DEA exemption, exemptions in 23 states, and the American Indian Religious Freedom Act).

\section{Procedural History}

The Oregon courts held the benefit denials unconstitutional under \textit{Sherbert v. Verner}. In \textit{Smith I} (1988), the Supreme Court vacated and remanded: if Oregon could constitutionally criminalize sacramental peyote use, it could deny benefits, but the state courts had not decided whether its criminal law reached such use. On remand, the Oregon Supreme Court held that the statute ``makes no exception for the sacramental use'' of peyote but that the ban violated the Free Exercise Clause. The Supreme Court again granted certiorari and reversed on April 17, 1990.

\section{Judicial Votes}

Reversed, \textbf{6--3} in the judgment, but \textbf{5--4} on the governing standard.
\begin{itemize}
\item \textbf{Majority:} Justice Scalia, joined by Chief Justice Rehnquist and Justices White, Stevens, and Kennedy.
\item \textbf{Concurring in the judgment:} Justice O'Connor; Justices Brennan, Marshall, and Blackmun joined Parts I and II, defending strict scrutiny.
\item \textbf{Dissent:} Justice Blackmun, joined by Justices Brennan and Marshall.
\end{itemize}

\section{Holding}

Because Oregon's criminal prohibition of peyote is constitutional, Oregon may deny unemployment benefits to employees dismissed for sacramental peyote use. The Free Exercise Clause does not relieve an individual of the obligation to comply with a ``valid and neutral law of general applicability.'' \textit{Sherbert}'s compelling-interest test does not apply to across-the-board, religion-neutral laws, and religious exemptions are left to the political process. The carve-outs for ``hybrid'' claims and for systems of individual exemptions explained prior precedent rather than decided this case, yet both later became battlegrounds in cases such as \textit{Fulton}.

\section{Analysis of Opinions}

\textbf{Majority (Scalia).} Government may not regulate belief or ban acts only when done for religious reasons, but an incidental burden from a neutral, generally applicable law does not ``prohibit'' free exercise. From \textit{Reynolds} through \textit{Lee}, the Court had never exempted religious conduct from such a law on free exercise grounds alone; \textit{Cantwell} and \textit{Yoder} were ``hybrid'' cases. \textit{Sherbert} had never invalidated anything outside unemployment compensation, whose ``good cause'' standard invites individualized assessment. Applying strict scrutiny to every religious objection would create a ``constitutional anomaly'' and invite exemptions from taxes, vaccination, and antidiscrimination laws, and judges cannot measure a practice's centrality. Minority faiths' relative disadvantage is an unavoidable consequence of democracy.

\textbf{Concurrence (O'Connor).} The text draws no line between targeted and general laws; the majority misread precedent and abandoned settled doctrine needlessly, and, as \textit{Barnette} teaches, the Bill of Rights withdraws minority liberties from politics. Writing alone, she found that Oregon's compelling interest in preventing harm and trafficking from Schedule I drugs required uniform enforcement, so the ban survived strict scrutiny.

\textbf{Dissent (Blackmun).} The majority effected ``a wholesale overturning of settled law.'' The relevant interest is refusing a religious exemption for ceremonial peyote, and it was merely symbolic: Oregon never prosecuted, offered no evidence of harm, faced negligible peyote trafficking, and the federal government and 23 states already allowed exemptions.

\section{Examples: Future Applications}

\begin{enumerate}
\item \textbf{Same-Side (No Strict Scrutiny): Vaccination Mandate.} A uniform school vaccination mandate with only medical exemptions is neutral and generally applicable; religious objectors must seek accommodation from the legislature.
\item \textbf{Same-Side (No Strict Scrutiny): Zoning.} A neutral ordinance barring all assembly halls, secular and religious, from a residential district gives a congregation no free exercise exemption, though RLUIPA may apply.
\item \textbf{Opposite-Side (Strict Scrutiny): Targeted Ritual Slaughter Ban.} An ordinance banning ``ritual'' animal sacrifice while exempting hunting and slaughterhouses is not neutral, as \textit{Lukumi} later held.
\item \textbf{Opposite-Side (Strict Scrutiny): Discretionary Exceptions.} A foster-care contract letting a commissioner grant exceptions in his ``sole discretion'' triggers strict scrutiny when an exception is refused to a religious agency, as in \textit{Fulton}.
\item \textbf{Fence-Sitter: Hybrid Claim.} A religious graphic designer challenges a neutral public-accommodations law by pairing free exercise with compelled speech; lower courts divide over whether the hybrid-rights dictum adds anything.
\end{enumerate}

\section{Critique}

\textbf{Backlash and doctrine.} A broad coalition persuaded Congress to pass RFRA in 1993; after \textit{City of Boerne} limited it, Congress passed RLUIPA and over twenty states enacted RFRAs. Scholars call the hybrid-rights distinction an unprincipled invention, find the individualized-exemptions carve-out unstable, and note that the parties never briefed overruling \textit{Sherbert}.

\textbf{Originalist and progressive views.} Michael McConnell argues that Founding-era sources support required exemptions, a view echoed by Justices Alito, Thomas, and Gorsuch in \textit{Fulton}; Philip Hamburger and Justice Barrett's caution defend \textit{Smith} or warn against overruling it without a replacement. Progressives, who once condemned \textit{Smith}, increasingly see it as a bulwark for equal citizenship and third-party rights.

\section{Key Quotations}

\begin{quote}
``We have never held that an individual's religious beliefs excuse him from compliance with an otherwise valid law prohibiting conduct that the State is free to regulate.''
\end{quote}
\noindent --- Justice Scalia, Opinion of the Court, 494 U.S. at 878--879.

\begin{quote}
``It may fairly be said that leaving accommodation to the political process will place at a relative disadvantage those religious practices that are not widely engaged in; but that unavoidable consequence of democratic government must be preferred to a system in which each conscience is a law unto itself or in which judges weigh the social importance of all laws against the centrality of all religious beliefs.''
\end{quote}
\noindent --- Justice Scalia, Opinion of the Court, 494 U.S. at 890.

\begin{quote}
``Precisely because `we are a cosmopolitan nation made up of people of almost every conceivable religious preference,' \textit{Braunfeld v. Brown}, 366 U.S. at 606, and precisely because we value and protect that religious divergence, we cannot afford the luxury of deeming presumptively invalid, as applied to the religious objector, every regulation of conduct that does not protect an interest of the highest order.''
\end{quote}
\noindent --- Justice Scalia, Opinion of the Court, 494 U.S. at 888.

\end{document}
